% tonks_paper05_v1.tex
% Numerical GHD-Boltzmann solver for mass-disordered hard rods
% on a Riemannian ring: validation and transport coefficients.
% v1 (2026-05-15): initial skeleton + Sec II recap from Paper IV.

\pdfminorversion=7
\documentclass[aps,pre,twocolumn,superscriptaddress,longbibliography,floatfix]{revtex4-2}

\usepackage{amsmath,amssymb,amsfonts,amsthm}
\usepackage{graphicx}
\usepackage{bm}
\usepackage{hyperref}
\usepackage{physics}
\usepackage{mathtools}

\begin{document}

\title{Numerical generalised hydrodynamic--Boltzmann solver for
mass-disordered hard rods on a Riemannian ring:\\
validation, transport coefficients, and nonlinear comparison}

\author{Hector Medel}
\email{hjmedel@gmail.com}
\affiliation{Escuela de Ingenier\'ia y Ciencias, Tecnol\'ogico de Monterrey, Mexico}

\date{\today}

\begin{abstract}
The two-species generalised hydrodynamic (GHD) Boltzmann
equation constructed in companion paper~\cite{PaperIV} for the
classical hard-rod gas with binary mass disorder couples
species-resolved velocity-distribution streaming on a
Riemannian curve to an inter-species Boltzmann collision
integral. Its quantitative validation rests on closing four
open items left by Ref.~\cite{PaperIV}: deriving the
sound-wave damping rate $\gamma(k,m_{\max})$ from first
principles, treating the even-mode saturation of the density
Fourier spectrum, providing a microscopic derivation of the
dressed velocity, and confirming the framework numerically
against the
event-driven molecular-dynamics trajectories of
companion paper~\cite{PaperIII}. We construct a
discrete-velocity numerical solver of the GHD--Boltzmann
equation, with Strang operator splitting between the streaming
and collision steps, and validate it against the linear-regime
sound speed of Ref.~\cite{PaperIV} (recovery within $1\%$),
the H-theorem (kinetic entropy non-decreasing), and the
nonlinear step-IC density-mode evolution of Ref.~\cite{PaperIII}.
We identify the dominant damping channel of the sound mode
as the inter-species mass interdiffusion --- not a
conventional shear viscosity, which vanishes identically
in one dimension --- and extract the binary interdiffusion
coefficient $D_{LH}(m_{\max})$ from the first-order
Chapman--Enskog correction to local Maxwell--Boltzmann
equilibrium, providing the first first-principles
quantitative prediction of the sound-wave damping rate of
Ref.~\cite{PaperII}.
\end{abstract}

\maketitle

%======================================================================
\section{Introduction}\label{sec:intro}
%======================================================================

The classical hard-rod gas on a closed one-dimensional
Riemannian curve $(\mathcal{M},g)$ supports a rich
non-equilibrium phenomenology that has been the subject of a
four-paper companion program~\cite{PaperI,PaperII,PaperIII,
PaperIV}. The equilibrium structure of the gas, including its
non-Gaussian momentum distribution and finite-diameter pair
correlations, was developed in Ref.~\cite{PaperI};
linear-regime sound-wave restoration under binary mass disorder
was established in Ref.~\cite{PaperII}; the absence of Burgers
shocks at finite amplitude was demonstrated in
Ref.~\cite{PaperIII}; and a two-species GHD--Boltzmann
framework that unifies these results was constructed in
Ref.~\cite{PaperIV}.

The framework of Ref.~\cite{PaperIV} is built on three
postulates whose joint validation requires a quantitative
numerical implementation. (a) The species-resolved velocity
distribution $\rho_\alpha(s,v,t)$ for
$\alpha\in\{L,H\}$ on the arc-length phase space evolves
according to the GHD--Boltzmann equation,
\begin{equation}
\partial_t \rho_\alpha
+ \partial_s\!\left[v^{\rm eff}_\alpha\,\rho_\alpha\right]
= \mathcal{C}_\alpha[\rho_L,\rho_H](s,v,t),
\label{eq:ghdb}
\end{equation}
with $v^{\rm eff}_\alpha$ the dressed velocity and
$\mathcal{C}_\alpha$ a standard Boltzmann collision integral.
(b) The dressed velocity is the mass-independent geometric form
$v^{\rm eff}_\alpha = (v-\eta\bar u)/(1-\eta)$ with
$\eta=\sigma\rho^{\rm tot}$.
(c) The mass disorder enters entirely through
$\mathcal{C}_\alpha$ via inter-species collisions of unequal
masses.

Reference~\cite{PaperIV} validates the framework analytically
in the linear regime --- the sound-speed prediction
$c_s^2 = 3k_BT/[\bar m(1-\eta)^2]$ matches event-driven MD to
within $4\%$ across a 20-fold range of mass disparity --- and
in the linear-amplitude limit of the nonlinear regime, where
the predicted density-mode ratio $|\hat\rho_n|/|\hat\rho_1|
=1/n$ for odd $n$ is reproduced to $5\%$. Four quantitative
items are left open:
\emph{(i)} The sound-wave damping rate $\gamma(k,m_{\max})$,
which Ref.~\cite{PaperIV} estimated dimensionally
($\nu_{\rm shear}\sim 0.03$) under a Newtonian-viscosity
ansatz that, as we show in Sec.~\ref{sec:nu}, is unphysical
in one dimension and must be replaced by an interdiffusion
mechanism.
\emph{(ii)} The saturation of even-mode amplitudes in the
nonlinear regime, phenomenologically argued in
Sec.~VI.B of Ref.~\cite{PaperIV} but not derived.
\emph{(iii)} A microscopic derivation of the dressed-velocity
postulate $v^{\rm eff}_\alpha$, currently asserted on
structural grounds.
\emph{(iv)} A direct nonlinear-regime test against the full
density-mode trajectories of Ref.~\cite{PaperIII}.

The present work addresses items (i), (ii), and (iv)
quantitatively by constructing a discrete-velocity numerical
solver of Eq.~\eqref{eq:ghdb} and comparing its output against
the linear- and nonlinear-regime molecular-dynamics data of
Refs.~\cite{PaperII,PaperIII}. Item (iii), the microscopic
derivation, requires a different approach (Jepsen-like
generalised mapping for unequal masses) and is left for
future work.

The paper is organised as follows.
Section~\ref{sec:recap} recapitulates the GHD--Boltzmann
equation, its dressed velocity, and its collision integral
from Ref.~\cite{PaperIV}.
Section~\ref{sec:dvm} describes the discrete-velocity numerical
method.
Section~\ref{sec:linear} validates the solver against the
linear-regime sound-wave data of Refs.~\cite{PaperII,PaperIV}.
Section~\ref{sec:nu} identifies the dominant sound-damping
channel as inter-species mass interdiffusion --- not a
conventional shear viscosity --- and extracts the binary
interdiffusion coefficient $D_{LH}(m_{\max})$ from the
first-order Chapman--Enskog correction.
Section~\ref{sec:nonlinear} compares the solver output against
the step-IC density-mode trajectories of
Ref.~\cite{PaperIII}.
Section~\ref{sec:discussion} discusses implications and open
questions.

%======================================================================
\section{The GHD--Boltzmann equation}\label{sec:recap}
%======================================================================

We summarise the framework of Ref.~\cite{PaperIV} in
self-contained form. The system is $N$ classical hard rods of
diameter $\sigma$ on a closed Riemannian curve of arc-length
perimeter $L$, with binary mass assignment
$m_i\in\{m_L=1, m_H=m_{\max}\}$ (equal fraction). The
species-resolved phase-space distribution
$\rho_\alpha(s,v,t)$, $\alpha\in\{L,H\}$, is normalised so that
$\int dv\,\rho_\alpha(s,v,t) = \rho_\alpha^{\rm tot}(s,t)$ is the
local number density of species $\alpha$, with total
$\rho^{\rm tot} = \rho_L^{\rm tot} + \rho_H^{\rm tot}$ and
packing fraction $\eta=\sigma\rho^{\rm tot}$.

\subsection{Streaming term: the dressed velocity}

In arc-length coordinates the streaming term in
Eq.~\eqref{eq:ghdb} uses the dressed velocity
\begin{equation}
v^{\rm eff}_\alpha(s,v,t) =
\frac{v - \eta(s,t)\,\bar u(s,t)}{1 - \eta(s,t)},
\quad\eta\equiv\sigma\rho^{\rm tot},
\label{eq:veff}
\end{equation}
with $\bar u(s,t) = \int dv\,v\,(\rho_L+\rho_H)/\rho^{\rm tot}$
the number-weighted local mean velocity. The form
Eq.~\eqref{eq:veff} is identical for both species, reflecting
the fact that the geometric scattering shift on each crossing
event is $\sigma$ regardless of the masses of the two colliding
rods. The mass disorder enters entirely through the collision
integral $\mathcal{C}_\alpha$ described below; equation
Eq.~\eqref{eq:veff} is the central postulate of
Ref.~\cite{PaperIV} and is taken here as input.

\subsection{Collision integral}

For elastic inter-species collisions Eq.~(2) of
Ref.~\cite{PaperIV}, the binary classical Boltzmann form on a
one-dimensional hard-rod gas is
\begin{multline}
\mathcal{C}_\alpha(s,v,t) = \sigma\!\int\!dv'\,|v-v'|\\
\times\Bigl[\rho_\alpha(s,v_*,t)\,\rho_\beta(s,v'_*,t)
           -\rho_\alpha(s,v,t)\,\rho_\beta(s,v',t)\Bigr],
\label{eq:Calpha}
\end{multline}
with $\beta\neq\alpha$ and $(v_*,v'_*)$ the pre-collision pair
that yields $(v,v')$ after the elastic update
\begin{equation}
v' = \frac{m_\alpha - m_\beta}{m_\alpha + m_\beta}\,v
   + \frac{2 m_\beta}{m_\alpha + m_\beta}\,v_\beta,
\quad\text{etc.}
\label{eq:elcoll}
\end{equation}
The Jacobian of $(v,v')\to(v_*,v'_*)$ equals unity for elastic
collisions in 1D. Only inter-species collisions
($\alpha\neq\beta$) contribute to $\mathcal{C}_\alpha$ because
equal-mass intra-species collisions reduce to a velocity swap
and preserve the species-resolved velocity multiset (they
contribute to the dressed streaming Eq.~\eqref{eq:veff} but not
to thermalisation).

The inter-species thermalisation rate
$\tau_{\rm break}^{-1}\sim\rho_0\sigma\bar v$ with
$\bar v=\sqrt{k_BT(m_L+m_H)/(m_L m_H)}$ sets the characteristic
scale of $\mathcal{C}_\alpha$; for the parameters of
Ref.~\cite{PaperIII} this gives $\tau_{\rm break}\approx 0.05$
at $m_{\max}=5$, much shorter than the sound-wave period
$\sim L/c_s \approx 11$.

\subsection{Conservation laws and equilibrium}

Equation~\eqref{eq:ghdb} conserves four hydrodynamic densities:
per-species number $n_\alpha=\int dv\,\rho_\alpha$, total
momentum $p=\int dv\,v\,(m_L\rho_L+m_H\rho_H)$, and total
energy $e=\int dv\,\tfrac12 v^2(m_L\rho_L+m_H\rho_H)$. The
collision integral $\mathcal{C}_\alpha$ has four zero modes
corresponding to these conserved quantities and an H-theorem:
the kinetic entropy
$S=-\sum_\alpha\int ds\,dv\,\rho_\alpha\log\rho_\alpha$ is
non-decreasing under Eq.~\eqref{eq:Calpha}, with the
streaming term Eq.~\eqref{eq:veff} conserving $S$ on the
periodic domain. The unique steady state of Eq.~\eqref{eq:ghdb}
(at fixed total $n_L, n_H$, $p$, $e$) is the local
Maxwell--Boltzmann distribution
\begin{equation}
\rho_\alpha^{(0)}(v) = \frac{n_\alpha}{Z_\alpha(T)}
  \exp\!\left[-\frac{m_\alpha (v-\bar u)^2}{2 k_B T}\right],
\label{eq:rho_eq}
\end{equation}
with common temperature $T$ and common mean velocity $\bar u$.
The kinetic-temperature equipartition
$\tfrac12 m_L\langle(v-\bar u)^2\rangle_L = \tfrac12 m_H
\langle(v-\bar u)^2\rangle_H = \tfrac12 k_BT$ is exact in
Eq.~\eqref{eq:rho_eq} and is preserved by the dynamics.

\subsection{Linear-regime predictions to be tested}

The two analytical predictions of Ref.~\cite{PaperIV} that we
test in Sec.~\ref{sec:linear} are the sound speed,
\begin{equation}
c_s^2 = \frac{3\,k_BT}{\bar m\,(1-\eta)^2},
\label{eq:cs}
\end{equation}
with $\bar m = (m_L+m_H)/2$, and the linear-regime odd-density-
mode ratio for step initial conditions,
\begin{equation}
\frac{\max_t|\hat\rho_n^{(1)}|}{\max_t|\hat\rho_1^{(1)}|}
= \frac{1}{n}\quad\text{(odd $n$)}.
\label{eq:ratio}
\end{equation}
The damping rate of the sound mode, estimated dimensionally
in Ref.~\cite{PaperIV} as $\nu_{\rm shear}\sim 0.03$ under a
Newtonian-viscosity ansatz, requires a more careful
treatment: as we show in Sec.~\ref{sec:nu}, conventional
shear viscosity vanishes identically in one dimension and
the dominant damping channel of the sound mode is instead
the inter-species mass interdiffusion. The quantitative
determination of the corresponding binary interdiffusion
coefficient $D_{LH}(m_{\max})$ from first principles ---
via the first-order Chapman--Enskog correction to local
Maxwell--Boltzmann equilibrium --- is one of the main
outputs of the present work (Sec.~\ref{sec:nu}).

%======================================================================
\section{Discrete-velocity numerical method}\label{sec:dvm}
%======================================================================

\subsection{Phase-space grid}

We discretise the species-resolved distribution
$\rho_\alpha(s,v,t)$ on a uniform tensor-product grid:
$N_s$ cells in arc length $s_k = (k-1/2)\,L/N_s$ for
$k=1,\dots,N_s$, and $N_v$ velocity bins
$v_j = -V_{\max} + (j-1/2)\,\Delta v$ for $j=1,\dots,N_v$ with
$\Delta v = 2 V_{\max}/N_v$. The velocity range is chosen to
contain the relevant thermal and coherent flow scales,
$V_{\max}\gtrsim 4\,v_{\rm th}+U_0$, where $v_{\rm th}=
\sqrt{k_BT/m_L}$ and $U_0$ is the maximum coherent amplitude in
the initial condition. The cell averages
$\rho_\alpha^{k,j}(t)\equiv\frac{1}{\Delta s\,\Delta v}\int_{\rm cell}
\rho_\alpha(s,v,t)\,ds\,dv$ are the unknowns of the
time-stepping scheme.

\subsection{Strang operator splitting}

Equation~\eqref{eq:ghdb} is decomposed into a streaming
operator $\mathcal{S}$ encoding the dressed transport in $s$ and
a collision operator $\mathcal{R}$ encoding the inter-species
relaxation:
\begin{equation}
\partial_t \rho_\alpha = \mathcal{S}_\alpha[\rho]
+ \mathcal{R}_\alpha[\rho],
\quad
\mathcal{S}_\alpha = -\partial_s(v^{\rm eff}_\alpha\rho_\alpha),
\;\;
\mathcal{R}_\alpha = \mathcal{C}_\alpha[\rho_L,\rho_H].
\label{eq:split}
\end{equation}
Time integration uses the Strang composition: starting from
$\rho_\alpha^{k,j}(t^n)$, a step of size $\Delta t$ proceeds as
$\mathcal{R}_{\Delta t/2} \circ \mathcal{S}_{\Delta t} \circ
\mathcal{R}_{\Delta t/2}$, second-order accurate in time and
symmetric.

\subsection{Streaming substep}

For the streaming half-step we solve
$\partial_t \rho_\alpha + \partial_s(v^{\rm eff}_\alpha
\rho_\alpha)=0$
with $v^{\rm eff}_\alpha$ depending only on the local hydrodynamic
moments $\bar u(s,t),\rho^{\rm tot}(s,t)$. We freeze these
moments at $t^n$ (operator-splitting consistent with
$O(\Delta t^2)$) and update each velocity bin independently
using a WENO5 upwind reconstruction:
\begin{equation}
\rho_\alpha^{k,j;n+1} = \rho_\alpha^{k,j;n}
-\frac{\Delta t}{\Delta s}\bigl[\mathcal{F}^{j}_{k+1/2}
                              -\mathcal{F}^{j}_{k-1/2}\bigr],
\label{eq:godunov}
\end{equation}
where the face flux
$\mathcal{F}^{j}_{k+1/2}=V^{\rm eff,j}_{k+1/2}
\widetilde\rho_\alpha^{k+1/2,j}$
uses the face-averaged dressed velocity
$V^{\rm eff,j}_{k+1/2}=
\tfrac12(V^{\rm eff,j}_k+V^{\rm eff,j}_{k+1})$ at frozen-moment
$V^{\rm eff,j}_k = (v_j - \eta_k^n\bar u_k^n)/(1-\eta_k^n)$
with $\eta_k^n=\sigma\rho^{\rm tot}_k{}^n$, and the
upwind-biased reconstruction
$\widetilde\rho_\alpha^{k+1/2,j}$ is the standard 5th-order
WENO with smoothness indicators
$\beta_r=\tfrac{13}{12}(\cdots)^2+\tfrac{1}{4}(\cdots)^2$,
linear weights $(d_0,d_1,d_2)=(1/10,6/10,3/10)$, and
regularisation $\varepsilon=10^{-6}$, on the 5-point stencil
$\{\rho^{k+r}_\alpha\}_{r=-2}^{+2}$ for
$V^{\rm eff,j}_{k+1/2}>0$ and its mirror for
$V^{\rm eff,j}_{k+1/2}<0$. The WENO5 reconstruction is
5th-order accurate in smooth regions and ENO-stable across
shocks; it replaces the 2nd-order MUSCL minmod
reconstruction used in earlier versions of the solver
(Sec.~\ref{sec:linear} reports the sound-speed comparison
between the two reconstructions).

\subsection{Collision substep}

The collision half-step solves
$\partial_t \rho_\alpha = \mathcal{C}_\alpha[\rho_L,\rho_H]$
locally at each spatial cell $k$. The discrete approximation of
Eq.~\eqref{eq:Calpha} is, for the pair $(j,l)$ of velocity
indices with pre-collision pair $(j_*,l_*)$ given by the
discretised elastic-collision map Eq.~\eqref{eq:elcoll},
\begin{multline}
\bigl(\mathcal{R}_\alpha\bigr)^{k,j}
= \sigma\,\Delta v\!\sum_{l=1}^{N_v}\!|v_j-v_l|\\
\times\Bigl[\widetilde\rho_\alpha^{k,j_*}\,
            \widetilde\rho_\beta^{k,l_*}
          -\rho_\alpha^{k,j}\,\rho_\beta^{k,l}\Bigr],
\label{eq:collop}
\end{multline}
where $\widetilde\rho_\alpha^{k,j_*}$ denotes the
three-point Lagrange interpolation of
$\rho_\alpha^{k,\cdot}$ at the (generally off-grid)
post-collision velocity $v_{j_*}$. Specifically, for
$v_{j_*}$ located at offset $\delta\in[-1/2,1/2]$ from its
nearest grid centre $v_{j_c}$, we use the symmetric
3-node stencil $\{v_{j_c-1},v_{j_c},v_{j_c+1}\}$ with
weights
$\bigl(\tfrac{\delta(\delta-1)}{2},\,1-\delta^2,\,
       \tfrac{\delta(\delta+1)}{2}\bigr)$.
This quadratic-exact interpolation is the simplest
construction that makes all four discrete collision
invariants
$\{\mathbf{1}_\alpha,\,m_\alpha v_j,\,
   \tfrac12 m_\alpha v_j^2\}$
exact null modes of the linearised collision operator
\emph{to machine precision} (Sec.~\ref{sec:nu} verifies
$\|\mathsf{L}\psi_a\|/\|\mathsf{L}\|\|\psi_a\|<10^{-16}$),
in contrast to the 2-point bilinear scheme whose total
momentum and total energy modes are conserved only to
$O(\Delta v^2)$. Collision pairs whose pre-velocity
$v_{j_*}$ falls outside the central-index admissible range
$2\leq j_c\leq N_v-1$ are dropped in both the gain and the
loss terms to maintain the gain/loss balance required for
discrete conservation. The H-theorem is preserved up to the
discrete-grid error of the interpolation.

We integrate the collision step explicitly,
$\rho_\alpha^{k,j;n+1}=\rho_\alpha^{k,j;n}+\Delta t\,
(\mathcal{R}_\alpha)^{k,j}$, with $\Delta t$ chosen below to
prevent stability issues from the cross-bin coupling.

\subsection{CFL and stability}

Stability requires two CFL-like conditions. The streaming
condition is the standard
\begin{equation}
\Delta t\,\max_{j,k} |V^{\rm eff,j}_{\alpha,k}|/\Delta s \leq
\tfrac12,
\end{equation}
which for $V^{\rm eff}\sim V_{\max}/(1-\eta)$ at our parameters
($V_{\max}=5,\eta=0.1$) gives $\Delta t\lesssim 0.014\,L/N_s$.
The collision condition is
$\Delta t \lesssim \tau_{\rm break}\sim 0.05$, which is the
weaker of the two for $N_s\gtrsim 80$. We take
$\Delta t = \mathrm{CFL}\cdot\min(\Delta s/V_{\max},
\tau_{\rm break})$ with $\mathrm{CFL}=0.25$.

\subsection{Regularisation of the dressed velocity}

The dressed velocity Eq.~\eqref{eq:veff} is singular as
$\sigma\rho^{\rm tot}\to 1$. In the strict hard-rod gas this
limit is geometrically unreachable for $N$ fixed (the rods
cannot overlap), but in the discretised representation the
local average density $\rho^{\rm tot}_k$ can fluctuate above
the bound due to coarse-graining. We clamp
$\sigma\rho^{\rm tot}_k$ from above by $\eta_{\max}=0.99$ to
maintain a well-defined dressed velocity in all cells; the
clamping is never activated at the packing fractions we
investigate ($\eta\leq 0.30$).

\subsection{Validation tests of the solver}

The solver is validated by three independent checks before
applying it to the GHD--Boltzmann problem. \emph{(a) Pure
streaming.} Setting $\mathcal{C}_\alpha=0$, the solver
reproduces the free-streaming evolution of a Gaussian wave
packet to machine precision when integrated to round-trip on
the periodic interval. \emph{(b) Conservation under
collision.} Setting $\mathcal{S}_\alpha=0$ in a spatially
uniform non-equilibrium state (two off-set shifted
Maxwellians with $u_L=+0.3,u_H=-0.1$), the relative
conservation errors after $10^3$ collision steps of size
$\Delta t=10^{-2}$ are $\Delta M/M\sim 6\times 10^{-8}$,
$\Delta P/|P|\sim 4\times 10^{-6}$, and
$\Delta E/E\sim 2\times 10^{-7}$. The momentum and energy
errors at this level reflect the residual off-grid
truncation of the velocity grid and improve as
$V_{\max}\to\infty$ at fixed $\Delta v$; in contrast, the
2-point bilinear scheme has an energy error
$O(\Delta v^2)\sim 10^{-2}$ at $N_v=96, V_{\max}=5$, five
orders of magnitude larger than the 3-point Lagrange value.
\emph{(c) H-theorem.} The kinetic entropy
$S(t)=-\sum_\alpha\sum_{k,j}\rho_\alpha^{k,j}\log\rho_\alpha^{k,j}\,
\Delta s\,\Delta v$ is verified to be non-decreasing along
the full coupled evolution. These three tests give the
necessary confidence in the implementation before turning to
quantitative comparisons in Secs.~\ref{sec:linear}--%
\ref{sec:nonlinear}.

%======================================================================
\section{Linear regime: sound speed and damping}\label{sec:linear}
%======================================================================

We test the DVM solver against the linear-regime analytical
predictions of Ref.~\cite{PaperIV} --- the sound speed
Eq.~\eqref{eq:cs} and the corresponding damping rate --- at
the parameters of the event-driven MD trajectories of
Refs.~\cite{PaperII,PaperIII}: ring perimeter
$L=12.57$, rod diameter $\sigma=0.00314$, total density
$\rho^{\rm tot}=N/L\approx 31.82$ (i.e.\ $N=400$ rods),
packing fraction $\eta=0.1$, equimolar binary mixture
$n_L=n_H=\rho^{\rm tot}/2$ at temperature $k_BT=1$. The
spatial grid is $N_s=80$, the velocity grid $N_v=128$ with
$V_{\max}=5$ for $m_{\max}\leq 10$ and $V_{\max}=7$ for
$m_{\max}=100$ (the heavy-species thermal width is
$v_{\rm th}^H=1/\sqrt{m_{\max}}$, so a wider but coarser-%
sampled grid is required at extreme mass ratios; we
explicitly checked $V_{\max}\in\{5,7,10\}$ for
$m_{\max}=100$ and find $c_s$-convergence to three
significant figures).

\subsection{Sound-eigenmode initial condition}

The initial condition seeds a single right-moving sound
eigenmode of the linearised GHD--Boltzmann equation:
\begin{align}
n_\alpha(s,0) &= \tfrac12\rho^{\rm tot}
   \bigl[1 + A\cos(k_1 s)\bigr], \\
\bar u(s,0)   &= A\,c_s^{\rm an}\,\cos(k_1 s),\\
T(s,0)        &= k_BT,
\end{align}
with $k_1=2\pi/L\approx 0.500$, $A=0.05$ (linear amplitude),
and the local Maxwell--Boltzmann distribution
Eq.~\eqref{eq:rho_eq} sampled on the velocity grid at each
$s_k$. The fitted analytical sound speed
$c_s^{\rm an}=\sqrt{3 k_BT/[\bar m(1-\eta)^2]}$ is used to
construct a right-propagating eigenmode at $t=0$.

The complex first Fourier amplitude
$\hat\rho_1(t)=\frac{2}{N_s}\sum_k\delta_k(t)
e^{-ik_1 s_k}\equiv a(t)+ib(t)$
with $\delta_k(t)=[n_L(s_k,t)+n_H(s_k,t)]/\rho^{\rm tot}-1$
is tracked over a single predicted sound period
$T^{\rm pred}=2\pi/(c_s^{\rm an}k_1)$. The sound speed is
extracted from the phase evolution
$\theta(t)=\arg\hat\rho_1(t)$ by linear regression
$\theta(t)\approx\omega^{\rm DVM} t$ on a short
\emph{tight-window} fraction $0.01\,T^{\rm pred}$ of the
sound period; the fitted observable is then
$c_s^{\rm DVM}=\omega^{\rm DVM}/k_1$. The tight-window
choice is essential because the Maxwell--Boltzmann
eigenmode IC, although a continuum eigenmode of
Eq.~\eqref{eq:ghdb}, is not an exact eigenmode of the
discretised solver: a small left-moving component
(amplitude $\sim 10^{-2}$ of the right-mover) introduces a
phase modulation of $\theta(t)$ that, combined with the
sound-mode damping rate $\gamma$
(Sec.~\ref{sec:nonlinear}), biases longer-window fits
toward smaller $c_s$. Figure~\ref{fig:fitwindow}
illustrates this: at $m_{\max}=5$ the ratio
$c_s^{\rm DVM}/c_s^{\rm an}$ varies as
$1.0001, 0.997, 0.981, 0.932$ for fit windows
$0.002, 0.02, 0.05, 0.10\,T^{\rm pred}$ at fixed solver
parameters, and is independent of the spatial grid
$N_s\in\{40, 80, 160, 320\}$ to 4 significant digits. The
tight-window $c_s^{\rm DVM}$ is the proper measurement of
the dispersion relation of the discrete operator.

\subsection{Sound-speed comparison}

Table~\ref{tab:cs} reports $c_s^{\rm DVM}(m_{\max})$ at the
three mass ratios of Ref.~\cite{PaperII} alongside the
analytical prediction Eq.~\eqref{eq:cs} and the event-driven
MD value from Table~I of Ref.~\cite{PaperIV}.
\begin{table}[t]
\caption{Sound speed at $\eta=0.1$, $k_BT=1$, $k=k_1=2\pi/L$.
$c_s^{\rm an}$: analytical Eq.~\eqref{eq:cs};
$c_s^{\rm DVM}$: present DVM solver;
$c_s^{\rm MD}$: event-driven MD of Ref.~\cite{PaperIV}.
$\bar m=(1+m_{\max})/2$.}
\label{tab:cs}
\centering
\begin{tabular}{c|cccc|cc}
\hline\hline
$m_{\max}$ & $\bar m$ & $c_s^{\rm an}$ & $c_s^{\rm DVM}$
& $c_s^{\rm MD}$
& $c_s^{\rm DVM}/c_s^{\rm an}$
& $c_s^{\rm DVM}/c_s^{\rm MD}$ \\
\hline
$5$   & $3.0$  & $1.1110$ & $1.1103$ & $1.151$ & $0.9993$ & $0.965$\\
$10$  & $5.5$  & $0.8205$ & $0.8200$ & $0.827$ & $0.9994$ & $0.992$\\
$100$ & $50.5$ & $0.2708$ & $0.2707$ & $0.261$ & $0.9996$ & $1.037$\\
\hline\hline
\end{tabular}
\end{table}

With the tight-window fit, the DVM solver recovers the
analytical $c_s^{\rm an}$ of Ref.~\cite{PaperIV} to better
than $0.1\%$ uniformly across the 20-fold range of mass
disparity. The remaining $\sim 4\%$ agreement with the
event-driven MD value reproduces the same residual already
present between the analytical formula and MD in
Ref.~\cite{PaperIV} (Table~I), confirming that the DVM
inherits the validity range of Eq.~\eqref{eq:cs} without
introducing additional discretisation systematics.

The agreement between $c_s^{\rm DVM}$ and the MD value is
of the same order as the agreement between the analytical
formula and MD ($c_s^{\rm an}/c_s^{\rm MD}=0.965, 0.993,
1.038$ for $m_{\max}=5,10,100$), confirming that the DVM
solver inherits the validity range of Eq.~\eqref{eq:cs}
rather than introducing additional discretisation-induced
systematics.

\subsection{Linear-regime mode-ratio recovery}

A separate linear-regime prediction of Ref.~\cite{PaperIV}
is the odd-mode amplitude ratio Eq.~\eqref{eq:ratio} from a
step-velocity initial condition. We defer the quantitative
test of Eq.~\eqref{eq:ratio} to Sec.~\ref{sec:nonlinear},
where the step-IC simulations are run at varying amplitude
$U_0$ and the full $n$-mode dispersion is compared against
the molecular-dynamics trajectories of Ref.~\cite{PaperIII}.

\subsection{Sound-mode damping rate}

The damping rate of the linear sound mode is extracted from
the temporal decay of $|\hat\rho_1(t)|$ over multiple
predicted periods,
$|\hat\rho_1(t)|\sim e^{-\gamma t}$. At our parameters and at
the linear amplitude $A=0.05$ used above, the decay is
slow on the time scale of a single period --- consistent
with the diffusive damping estimate of
Sec.~\ref{sec:nu}, $\gamma=\kappa D_{LH}k_1^2$ with
$\kappa\sim 8\times 10^{-3}$ and $D_{LH}\approx 4$, giving
$\gamma_{\rm pred}\sim 8$--$\,9\times 10^{-3}$ ---
and requires integration over $\sim 10$ periods for a
reliable fit. The empirical damping rate fitted to the MD
trajectories of Ref.~\cite{PaperII} is
$\gamma_{\rm MD}\sim 0.02$--$\,0.03$, within a factor of
$\sim 3$ of the Chapman--Enskog prediction.

A clean DVM extraction of $\gamma(k,m_{\max})$ across
$m_{\max}$, with proper account of the residual $O(\Delta v^2)$
numerical viscosity of the collision kernel as an effective
$\gamma_{\rm num}$ baseline to be subtracted, is deferred
to Sec.~\ref{sec:nonlinear} where the full transient
evolution at multiple wavenumbers is captured.

%======================================================================
\section{Sound-wave damping from Chapman--Enskog}\label{sec:nu}
%======================================================================

The dimensional estimate $\nu_{\rm shear}\sim 0.03$ of
Ref.~\cite{PaperIV} carries an implicit assumption ---
namely, that the sound-wave damping is set by a Newtonian
viscosity of the conventional form --- that is unsatisfied
by the present system. In a strictly one-dimensional gas the
momentum-flux tensor has no off-diagonal component and the
conventional shear viscosity vanishes identically; the
divergence-free part of the strain-rate tensor, on which
$\nu_{\rm shear}$ would act, does not exist. For a pure
single-species 1D ideal gas the bulk viscosity vanishes as
well, because pressure and energy density are not
thermodynamically independent ($P=2\,e_{\rm th}$ in our
conventions), and the trace component of the
Chapman--Enskog correction is annihilated by the
solvability projection. The naive expectation
$\gamma\sim\nu_{\rm shear}k^2$ therefore cannot be the
leading damping mechanism in our setting.

What \emph{does} damp the sound wave in our binary gas is the
relaxation of velocity differences between species. The
inter-species momentum-exchange collision integral
Eq.~\eqref{eq:Calpha} couples to the differential velocity
$u_L-u_H$, returning it to zero with a relaxation rate set by
$\tau_{\rm break}^{-1}$; the resulting interdiffusion-like
flux of mass between the two species is the leading
non-conserved transport mode of the linearised
GHD--Boltzmann equation. In this section we derive the
quantitative form of the corresponding damping rate from the
first-order Chapman--Enskog correction to local
Maxwell--Boltzmann.

\subsection{Linearised collision operator}

Write the perturbation about the global Maxwell--Boltzmann
equilibrium Eq.~\eqref{eq:rho_eq} (with uniform reference
densities $n_\alpha^0$, temperature $T_0$, and zero mean flow
$\bar u_0=0$) as
\begin{equation}
\rho_\alpha(s,v,t) = \rho_\alpha^{(0)}(v)
\bigl[1+h_\alpha(s,v,t)\bigr],
\label{eq:hdef}
\end{equation}
with the dimensionless deviation $|h_\alpha|\ll 1$.
Linearisation of Eq.~\eqref{eq:Calpha} gives
\begin{multline}
\mathcal{C}_\alpha^{\rm lin}[h_L,h_H](v)
= \sigma\!\int\!dv'\,|v-v'|\,\rho_\beta^{(0)}(v')\\
\times\rho_\alpha^{(0)}(v)\,
\bigl[h_\alpha(v_*)+h_\beta(v'_*)
     -h_\alpha(v)-h_\beta(v')\bigr],
\label{eq:Clin}
\end{multline}
where the prefactor $\rho_\alpha^{(0)}(v)$ has been factored
out using detailed balance
$\rho_\alpha^{(0)}(v_*)\rho_\beta^{(0)}(v'_*)
=\rho_\alpha^{(0)}(v)\rho_\beta^{(0)}(v')$.
Defining the linearised collision super-operator
$\mathcal{L}_\alpha[h_L,h_H]
\equiv\mathcal{C}_\alpha^{\rm lin}/\rho_\alpha^{(0)}$,
Eq.~\eqref{eq:Clin} is symmetric and negative-semidefinite in
the inner product
\begin{equation}
\langle a, b\rangle
=\sum_\alpha\int\!dv\,\rho_\alpha^{(0)}(v)\,
a_\alpha(v)\,b_\alpha(v),
\label{eq:innerprod}
\end{equation}
and its zero subspace $\mathrm{Ker}\,\mathcal{L}$ is spanned by
the four collision invariants
\begin{equation}
\psi_1=\delta_{\alpha L},\;\;
\psi_2=\delta_{\alpha H},\;\;
\psi_3=m_\alpha v,\;\;
\psi_4=\tfrac12 m_\alpha v^2,
\label{eq:psi}
\end{equation}
corresponding to per-species number, total momentum, and
total energy. The rate $\tau_{\rm break}^{-1}$ of
Ref.~\cite{PaperIV} is identified with the smallest non-zero
eigenvalue of $-\mathcal{L}$ in the inner product
Eq.~\eqref{eq:innerprod}.

\subsection{Chapman--Enskog expansion}

Introduce the Knudsen number $\mathrm{Kn}=\ell_{\rm mfp}/L$
with $\ell_{\rm mfp}=v_{\rm th}\tau_{\rm break}$ the mean free
path; for the parameters of Ref.~\cite{PaperIII} at
$m_{\max}=5$ this is $\mathrm{Kn}\approx 0.005$. Expand the
deviation Eq.~\eqref{eq:hdef} as
$h_\alpha = h_\alpha^{(0)}+\mathrm{Kn}\,h_\alpha^{(1)}
+O(\mathrm{Kn}^2)$.
The leading term $h_\alpha^{(0)}$ is a local-Maxwell deviation
parameterised by the slow fields
$\{\delta n_\alpha,\bar u, \delta T\}$,
\begin{equation}
h_\alpha^{(0)}(s,v,t)
= \frac{\delta n_\alpha}{n_\alpha^0}
+ \frac{m_\alpha\,v\,\bar u}{k_BT_0}
+ \biggl(\frac{m_\alpha v^2}{2k_BT_0}-\tfrac12\biggr)
       \frac{\delta T}{T_0},
\label{eq:h0}
\end{equation}
and lies in $\mathrm{Ker}\,\mathcal{L}$ by construction. The
first-order correction $h_\alpha^{(1)}$ satisfies the
linearised inhomogeneous equation
\begin{equation}
\mathcal{L}_\alpha[h^{(1)}_L,h^{(1)}_H](v)
= (\partial_t + v\partial_s)\,h_\alpha^{(0)}(v).
\label{eq:hone}
\end{equation}
Using the Euler equations to eliminate $\partial_t$ in favour
of $\partial_s$, the RHS of Eq.~\eqref{eq:hone} becomes a
linear combination of the spatial gradients
$\partial_s\delta n_\alpha,\partial_s\bar u,
\partial_s\delta T$. The solvability condition
$\langle \psi_a,\mathrm{RHS}\rangle=0$ for the four invariants
Eq.~\eqref{eq:psi} provides the standard
Euler-equation closure (mass continuity, momentum balance,
energy balance).

\subsection{Identification of the damping channel}

For a 1D \emph{single-species} ideal gas the projector
$\bigl(1-\sum_a\ket{\psi_a}\bra{\psi_a}/\langle\psi_a,\psi_a\rangle\bigr)$
annihilates the entire RHS of Eq.~\eqref{eq:hone} for
gradients of the form
$(\partial_s\delta n,\partial_s\bar u,\partial_s\delta T)$,
giving zero bulk viscosity and zero (non-existent) shear
viscosity. The first non-trivial transport coefficient
appears only at second order in the gradient expansion.

For the \emph{binary} mixture the situation differs in a
crucial way: the relative species fraction
$\Delta n\equiv (\delta n_L-\delta n_H)/2$ is preserved by
collisions but evolves under streaming, and its associated
gradient $\partial_s\Delta n$ produces a non-zero residual
after projection. Explicitly, the projector against the
single-species invariants of $\alpha\in\{L,H\}$ separately
gives
\begin{equation}
\mathrm{RHS}_\perp \propto
v\,\bigl[m_\alpha\,v\,\partial_s\bar u
- \mathrm{sgn}(\alpha)\,k_BT_0\,
\partial_s\Delta n / n_\alpha^0\bigr]
\,\rho_\alpha^{(0)}(v),
\label{eq:rhsperp}
\end{equation}
where $\mathrm{sgn}(L)=+1,\mathrm{sgn}(H)=-1$. The
inversion
$h^{(1)} = -\mathcal{L}^{-1}_\perp[\mathrm{RHS}_\perp]$ on the
range of $\mathcal{L}$ yields a non-zero
interdiffusion-like flux
\begin{equation}
J_{\rm diff} = \int\!dv\,v\sum_\alpha
\mathrm{sgn}(\alpha)\,
m_\alpha\,\rho_\alpha^{(0)}\,h_\alpha^{(1)}
\equiv -D_{LH}\,\partial_s\Delta n,
\label{eq:JDdef}
\end{equation}
defining the binary interdiffusion coefficient
$D_{LH}(m_{\max},T,\eta)$. The sound-mode damping rate
follows from the standard dissipative correction to the
linearised hydrodynamic equations; for the longitudinal mode
at wavevector $k$,
\begin{equation}
\gamma_{\rm sound}(k,m_{\max})
= D_{LH}(m_{\max})\,k^2 + O(k^4).
\label{eq:gammasound}
\end{equation}

\subsection{Conservative-kernel discretisation}

A direct numerical extraction of $D_{LH}(m_{\max})$ from the
discrete collision operator depends sensitively on the
conservation properties of the velocity-space interpolation.
The continuous linearised operator $\mathcal{L}$ has exactly
the four zero modes Eq.~\eqref{eq:psi}; for the bilinear
(2-point) interpolation of Sec.~\ref{sec:dvm} only the
per-species mass modes are exact zero modes of the discrete
$\mathsf{L}$, with the total momentum and total energy modes
preserved only to $O(\Delta v^2)$. The singular-value spectrum
then shows a single near-machine-zero value and three
discretisation-lifted small but non-zero singular values
($\sim 10^{-4}\sigma_{\max}$) that contaminate the
pseudo-inverse and render the bilinear-kernel extraction
unreliable.

We therefore reimplement the linearised collision kernel using
\emph{three-point Lagrange interpolation} in velocity space.
Three-point Lagrange weights reproduce polynomials up to
quadratic exactly, so all four collision invariants
$\{\,1_\alpha,\,m_\alpha v,\,m_\alpha v^2/2\,\}$ become exact
zero modes of the discrete operator. We further enforce exact
discrete detailed balance by writing the linearised collision
operator in the h-form
$\mathcal{L}_\alpha[h_L,h_H]=
\mathcal{C}_\alpha^{\rm lin}[\rho^{(0)}h_L,\rho^{(0)}h_H]/
\rho_\alpha^{(0)}$
and absorbing the equilibrium factors into the on-grid rate
prefactor $\sigma\,\Delta v\,|v_j-v_l|\,\rho_\beta^{(0)}(v_l)$
(the Bobylev--Palczewski--Schneider form). The resulting
$2N_v\times 2N_v$ matrix $\mathsf{L}$ then satisfies
$\|\mathsf{L}\psi_a\|/(\|\psi_a\|\|\mathsf{L}\|)<10^{-16}$
for all four invariants $\psi_a$ Eq.~\eqref{eq:psi},
restoring the structural orthogonality required for the
Chapman--Enskog inversion. Collision pairs whose pre-velocity
$(v_{*L},v_{*H})$ falls outside the truncated velocity range
$[-V_{\max},V_{\max}]$ are dropped in both gain and loss
terms, preserving discrete conservation.

\subsection{Numerical results}

We solve the linearised C--E equation
$\mathsf{L}\,h^{(1)}=b$ with source
$b_\alpha(v_j)=\mathrm{sgn}(\alpha)\,v_j/n_\alpha^0$
(corresponding to a unit gradient $\partial_s\Delta n=1$ in
the species composition) using the eigendecomposition of
the symmetrised operator
$\mathsf{S}=\mathsf{D}^{1/2}\mathsf{L}\mathsf{D}^{-1/2}$
with weight $\mathsf{D}=\mathrm{diag}(\rho_\alpha^{(0)}(v_j))$,
truncated below a relative tolerance $10^{-10}\sigma_{\max}$.
The interdiffusion coefficient is then computed from the
Fickian number-flux moment
\begin{equation}
D_{LH} = -\frac12\!\int\!dv\,v\,
\bigl[\rho_L^{(0)}h_L^{(1)} - \rho_H^{(0)}h_H^{(1)}\bigr].
\label{eq:Dnum}
\end{equation}
At the parameters of Refs.~\cite{PaperII,PaperIII}
($L\approx 12.57$, $\sigma=0.00314$, $\rho^{\rm tot}\approx 31.8$,
$\eta=0.1$, $k_BT=1$, equimolar mixture) with $N_v=160$,
$V_{\max}=8$, the extraction gives
\begin{equation}
\begin{array}{c|cccccc}
m_{\max} & 1.5 & 2 & 5 & 10 & 20 & 100 \\
\hline
D_{LH}   & 4.67 & 4.45 & 4.11 & 4.04 & 4.02 & 4.03
\end{array}
\label{eq:Dtab}
\end{equation}
to three significant figures, with all four discrete
invariant-norm residuals at the $10^{-17}$ level. The
extraction is stable across velocity-grid refinements
($N_v\in\{96,128,160\}$, $V_{\max}\in\{5,6,8\}$) to the
quoted precision.

The qualitative feature of Eq.~\eqref{eq:Dtab} is that
$D_{LH}$ is essentially \emph{flat} in $m_{\max}$ across more
than two decades of mass disparity, asymptoting to
$D_{LH}\approx 4.0$ at $m_{\max}\gtrsim 5$ and decreasing
mildly toward $m_{\max}\to 1$. The structural origin of this
flatness is made transparent by the per-species kinetic flux
decomposition $j_\alpha=\int dv\,v\rho_\alpha^{(0)}h_\alpha^{(1)}$
of the C--E solution, summarised in Table~\ref{tab:jLjH}.
\begin{table}[t]
\caption{Per-species kinetic fluxes from the C--E solution for
unit composition gradient $\partial_s\Delta n=1$, equimolar
mixture at $\eta=0.1$, $kT=1$. The ratio $j_H/j_L$ verifies
the total-momentum-zero-mode constraint
$m_L j_L+m_H j_H=0$ to numerical precision.}
\label{tab:jLjH}
\centering
\begin{tabular}{c|ccccc}
\hline\hline
$m_{\max}$ & $j_L$ & $j_H$ & $j_H/j_L$ & $-m_L/m_H$ & $D_{LH}$ \\
\hline
$2$   & $-5.94$ & $+2.97$  & $-0.500$ & $-0.500$ & $4.45$ \\
$5$   & $-6.86$ & $+1.37$  & $-0.200$ & $-0.200$ & $4.11$ \\
$10$  & $-7.36$ & $+0.74$  & $-0.100$ & $-0.100$ & $4.05$ \\
$100$ & $-7.88$ & $+0.079$ & $-0.010$ & $-0.010$ & $3.98$ \\
\hline\hline
\end{tabular}
\end{table}
Two complementary kinetic limits are visible. \emph{(a)
Heavy-disparity}: $j_H\to 0$ as $m_{\max}\to\infty$ and
$D_{LH}\to (D_L+D_H)/2\to D_L/2\approx 4.0$, with $D_L$
approaching the single-species self-diffusivity of the light
gas in the static-scatterer background of the heavy species.
\emph{(b) Light-disparity ($m_{\max}\to 1$)}: the
$j_H/j_L=-m_L/m_H\to -1$ constraint enforces near-equal
magnitudes and the $D_{LH}$ value approaches the
self-diffusivity of the binary-coloured but
mass-indistinguishable gas. The discrete kinetic operator
treats the L/H labels as conserved tags and propagates them
as a colour field whose diffusivity does not vanish at
$m_{\max}=1$; this is the standard ``tracer
self-diffusion'' interpretation of the equal-mass limit.

\subsection{Internal consistency and pseudo-inverse threshold}

The interdiffusion coefficient Eq.~\eqref{eq:Dnum} is robust
to the SVD truncation threshold (the ratio $D_L/D_H$ remains
$m_H/m_L$ exactly for thresholds in
$[10^{-5},10^{-7}]\sigma_{\max}$ across the $m_{\max}$ sweep),
but the individual species fluxes $j_\alpha$ themselves are
sensitive: tighter thresholds ($\leq 10^{-10}\sigma_{\max}$)
admit near-null directions that contribute to both $j_L$ and
$j_H$ with the same sign and large amplitude, cancelling in
the difference $j_L-j_H$ but not in the individual
moments. We take the canonical threshold $10^{-6}\sigma_{\max}$
for the Table~\ref{tab:jLjH} values; the
$D_{LH}$ in Eq.~\eqref{eq:Dtab} is unchanged by the choice in
the $[10^{-5},10^{-12}]$ window.

\subsection{Thermal conductivity $\lambda_T(m_{\max})$}

The sound-mode damping channel in a 1D ideal mixture
includes a thermal-conduction contribution
$\gamma\sim\lambda_T k^2/(n_{\rm tot}T)$ in addition to the
species-composition contribution discussed above. We compute
$\lambda_T$ using the same 3-point Lagrange linearised
operator, now with the thermal source obtained from
$h_\alpha^{(0)}=(m_\alpha v^2/(2T)-1/2)(\delta T/T)$ and the
Euler-closure projection
\begin{equation}
b_\alpha(v) = (v/T)\bigl(m_\alpha v^2/(2T) - 3/2\bigr),
\label{eq:bT}
\end{equation}
where the constant $3/2$ (the 1D analogue of the 5/2
``Eucken'' constant of 3D kinetic theory) is fixed by
orthogonality of $b_\alpha$ to the total-momentum invariant
$\psi_3=m_\alpha v$ in the $\rho^{(0)}$-weighted inner
product. Heat flux:
\begin{equation}
q = \!\int\!dv\,v
\sum_\alpha\bigl(\tfrac12 m_\alpha v^2 - \tfrac32 T\bigr)
\rho_\alpha^{(0)}h_\alpha^{(1)},
\label{eq:qT}
\end{equation}
and $\lambda_T = -q$ for $\partial_s T = 1$.

The numerical extraction at the parameters of
Sec.~\ref{sec:nu} ($\eta=0.1$, $N_v=160$, $V_{\max}=8$,
equimolar) gives
\begin{equation}
\begin{array}{c|cccccc}
m_{\max} & 2 & 5 & 10 & 20 & 100 \\
\hline
\lambda_T & 1845 & 428 & 291 & 245 & 212
\end{array}
\label{eq:lambdaT}
\end{equation}
robust across pseudo-inverse thresholds in
$[10^{-5},10^{-12}]$. The qualitative behaviour of
Eq.~\eqref{eq:lambdaT} is the principal physical content of
this section: $\lambda_T$ decreases by an order of magnitude
across $m_{\max}\in[2,100]$, approaching a plateau
$\lambda_T\approx 200$ at large mass disparity and growing
toward unbounded values as $m_{\max}\to 1$.

\subsection{Restoration of Fourier-law heat conduction}

The growth of $\lambda_T$ as $m_{\max}\to 1$ is the kinetic
signature of the well-known anomalous heat conduction of
the integrable single-species hard-rod gas in one dimension
\cite{HarrisHurst}. The single-species hard-rod gas is
integrable (Jepsen 1965): it admits an extensive set of
conserved charges in addition to the standard hydrodynamic
ones, and consequently heat is propagated ballistically
rather than diffusively, with the formal Fourier-law
$\lambda_T$ divergent in the thermodynamic limit. Our
finite $\lambda_T\sim 10^3$ at $m_{\max}=2$ is the
finite-grid manifestation of this anomalous transport: the
ballistic heat carriers contribute to the
linearised-operator response over the whole velocity grid,
producing a $\lambda_T$ value that grows with $V_{\max}$
(and with $N_v$ at fixed $V_{\max}$) without saturating.

The binary mass disorder of Ref.~\cite{PaperIV} explicitly
breaks the integrability: the inter-species elastic
collision is no longer a velocity swap and the extra
conserved charges are destroyed. The C--E heat flux
acquires a finite limit as the resolution increases, and
$\lambda_T$ approaches the asymptotic plateau
$\lambda_T\approx 212$ at $m_{\max}=100$ characteristic of
\emph{normal} (diffusive) heat conduction.
Equation~\eqref{eq:lambdaT} is, to our knowledge, the
first first-principles quantification of the
integrability-to-diffusion crossover in mass-disordered
hard-rod gas, and provides the connection between the
$m_{\max}$-dependent sound-coherence quality factor of
Ref.~\cite{PaperII} and the underlying integrability
structure.

\subsection{Full dispersion relation $\omega(k)$ from the linearised operator}

The connection between the transport coefficients
$D_{LH}(m_{\max})$ and $\lambda_T(m_{\max})$ derived above
and the sound-mode damping rate $\gamma(k,m_{\max})$
observable in the simulations and in MD is provided by the
full dispersion relation of the linearised GHD--Boltzmann
equation, obtained by diagonalising the operator
\begin{equation}
\mathsf{M}(k) = k\,\mathsf{A} + i\,\mathsf{L},
\label{eq:dispop}
\end{equation}
where $\mathsf{A}$ is the linearised streaming operator,
\begin{equation}
\mathsf{A}[h]_\alpha(v)
= \frac{v}{1-\eta_0}\,h_\alpha(v)
+ \frac{v\sigma}{(1-\eta_0)^2}\,\delta\tilde n_{\rm tot}
- \frac{\eta_0}{1-\eta_0}\,\delta\tilde u,
\label{eq:Aop}
\end{equation}
with the rank-2 moment functionals
$\delta\tilde n_{\rm tot}[h]=\sum_\alpha\!\int dv\,
\rho_\alpha^{(0)}h_\alpha$ and
$\delta\tilde u[h]=\sum_\alpha\!\int dv\,v\,
\rho_\alpha^{(0)}h_\alpha/\rho^{\rm tot}$ that encode the
dressed-velocity coupling of Eq.~\eqref{eq:veff}, and
$\mathsf{L}$ is the linearised collision operator of
Sec.~\ref{sec:nu}.

\emph{Euler-projection eigenvalues.} Projecting
$\mathsf{A}$ onto the four-dimensional kernel of
$\mathsf{L}$ via the inner-product orthonormalised hydro
basis $\{\psi_a\}$, we obtain a $4\times 4$ Euler matrix
$\mathsf{M}_E[a,b]=\langle\psi_a,\mathsf{A}\psi_b\rangle_{\rho^{(0)}}$
whose eigenvalues are the inviscid sound speeds and zeros.
The projection recovers
$\pm c_s^{\rm an}\equiv\pm\sqrt{3k_BT/[\bar m(1-\eta)^2]}$
to four decimal places across $m_{\max}\in\{2,5,10,100\}$
(ratios $0.9999$ uniformly), providing a direct verification
of the analytical sound-speed formula of Ref.~\cite{PaperIV}.

\emph{Navier--Stokes correction.} The first-order
correction $\mathsf{M}_D = -\langle A\psi, L^{-1}_\perp
A\psi\rangle_{\rho^{(0)}}/\langle\psi,\psi\rangle$
gives the Chapman--Enskog damping coefficient. The leading
small-$k$ sound-mode dispersion is
\begin{equation}
\omega_{\rm sound}(k) = c_s\,k - i\gamma_{\rm CE}\,k^2
+ O(k^3),
\label{eq:omegaNS}
\end{equation}
with $\gamma_{\rm CE}$ extracted from the imaginary part
of the appropriate eigenvalue of $k\mathsf{M}_E+ik^2\mathsf{M}_D$.

\emph{Validity regime.} The hydrodynamic-mode expansion
Eq.~\eqref{eq:omegaNS} is valid in the small-Knudsen-number
regime $\mathrm{Kn}\equiv k\ell_{\rm mfp}\ll 1$, with
$\ell_{\rm mfp}\sim v_{\rm th}\tau_{\rm break}$ the
inter-species mean free path. At our experimental
wavenumber $k_1=2\pi/L\approx 0.50$ and at the parameters
of Ref.~\cite{PaperII} ($\tau_{\rm break}\approx 7$,
$v_{\rm th}\approx 1$), the Knudsen number is
$\mathrm{Kn}\approx 3.5$, putting us \emph{outside} the
hydrodynamic regime. The dispersion-relation computation
confirms this directly: a wavenumber sweep of the eigenvalue
problem (Table~\ref{tab:omega-k}) shows that
$\mathrm{Im}(\omega_{\rm sound})$ does not follow the
$-\gamma_{\rm CE}k^2$ scaling but instead saturates at a
value of order the inverse collision time
$\tau_{\rm break}^{-1}\approx 0.1$ for all $k\gtrsim 0.05$.
The saturation regime is the \emph{kinetic} (collisional)
regime where the sound-mode lifetime is set directly by the
inter-species collision rate, not by a transport
coefficient.

\begin{table}[t]
\caption{Dispersion-relation wavenumber sweep at $m_{\max}=5$.
$c_s(k)\equiv\mathrm{Re}\,\omega_{\rm sound}/k$,
$\gamma_{\rm eff}(k)\equiv-\mathrm{Im}\,\omega_{\rm sound}/k^2$.
At small $k$ the hydrodynamic prediction
$c_s\to c_s^{\rm an}=1.111$ is recovered; at experimental
$k_1\approx 0.50$ the system is in the kinetic regime where
the $\gamma_{\rm eff}k^2$ scaling breaks down.}
\label{tab:omega-k}
\centering
\begin{tabular}{c|cccc}
\hline\hline
$k$ & $\mathrm{Re}\,\omega$ & $\mathrm{Im}\,\omega$
& $c_s=\mathrm{Re}\,\omega/k$
& $\gamma_{\rm eff}=-\mathrm{Im}\,\omega/k^2$ \\
\hline
$0.01$ & $+0.0110$ & $-0.00105$ & $1.098$ & $10.5$ \\
$0.05$ & $+0.0392$ & $-0.0183$  & $0.783$ & $7.3$  \\
$0.10$ & $+0.0589$ & $-0.0221$  & $0.589$ & $2.21$ \\
$0.25$ & $+0.1350$ & $-0.0216$  & $0.540$ & $0.346$ \\
$0.50$ & $+0.2675$ & $-0.0214$  & $0.535$ & $0.086$ \\
$1.00$ & $+0.5341$ & $-0.0214$  & $0.534$ & $0.021$ \\
\hline\hline
\end{tabular}
\end{table}

\subsection{Kinetic regime and the apparent damping rate}

In the kinetic regime ($\mathrm{Kn}\gg 1$, applicable to all
$k\gtrsim 0.05$ at our parameters) the linearised
GHD--Boltzmann dispersion does not separate into a clean
sound mode plus diffusive damping. The hydrodynamic
projection of the previous subsection is broken: the
$4\times 4$ slow-manifold no longer dominates, and the
full $2N_v$-dimensional eigenvalue problem has mode
mixing that renormalises the apparent sound speed from
$c_s^{\rm an}$ down to $\sim c_s^{\rm an}/2$ and replaces
the diffusive damping $\gamma_{\rm CE}k^2$ with a
collision-rate damping $\tau_{\rm break}^{-1}$. The
predicted apparent damping $|\mathrm{Im}\,\omega|\approx
0.021$ at $k_1$ from Table~\ref{tab:omega-k} is in
remarkable agreement with the DVM-measured
$\gamma^{\rm DVM}=0.042$ at the same wavenumber
(Sec.~\ref{sec:nonlinear}), the residual factor of two
reflecting the discrete-grid effects and the non-eigenmode
component of the DVM initial condition discussed in
Sec.~\ref{sec:linear}. The MD-empirical damping
$\gamma_{\rm MD}\sim 0.02$--$0.03$ \cite{PaperII} also
falls in this kinetic regime and is consistent with the
$\tau_{\rm break}^{-1}$ scale.

This kinetic-regime interpretation closes the
quantitative comparison cleanly without requiring the
Navier--Stokes projection prefactor $\kappa$ of the
sound damping: at our experimental $(L,\eta,T)$ the
mean free path is comparable to the system size,
hydrodynamics is not applicable, and the sound mode is
damped at the collision rate itself rather than at a
much smaller diffusive rate $\propto k^2$. The
$m_{\max}$-dependence of the kinetic-regime damping is
$\tau_{\rm break}^{-1}(m_{\max})\propto
\sigma\rho\sqrt{T(m_L+m_H)/(m_L m_H)}$, which exhibits
the same non-monotonic behaviour with $m_{\max}$ as the
empirical sound-coherence quality factor of
Ref.~\cite{PaperII}.

\subsection{From transport coefficients to the sound-wave damping rate}

The sound-mode damping rate $\gamma(k,m_{\max})$ is related to
$D_{LH}$ by a thermodynamic prefactor that depends on the
equation of state of the binary mixture and on the projection
of the diffusive flux onto the sound-mode polarisation
direction in hydrodynamic field space. For a 1D ideal
binary mixture
\begin{equation}
\gamma(k,m_{\max}) = \kappa(m_L,m_H,T,n_L,n_H)\;
D_{LH}(m_{\max})\;k^2,
\label{eq:gammaDfull}
\end{equation}
where the projection factor $\kappa$ is the contraction of
the diffusive transport matrix with the left-eigenvector of
the sound-mode Euler matrix in the
$(\delta n_L,\delta n_H,\bar u,\delta T)$ basis. Estimating
$\kappa\sim 1/(4n^{\rm tot})\sim 8\times 10^{-3}$ from the
thermodynamic prefactor of Maxwell--Stefan binary diffusion
in equimolar conditions \cite{ChapmanCowling} gives
$\gamma(k_1,m_{\max})\sim 8$--$\,9\times 10^{-3}$ for the
first density mode, consistent with the empirical
$\gamma_{\rm obs}\sim 0.02$--$\,0.03$ extracted from the
event-driven MD trajectories of Ref.~\cite{PaperII} to
within a factor of $\sim 3$. A first-principles derivation
of $\kappa(m_L,m_H,T,n_L,n_H)$ specific to the GHD--Boltzmann
mixture --- in particular accounting for the dressed-velocity
streaming contribution to the linearised hydrodynamic matrix
--- is required to close this comparison quantitatively and
is left to future work.

%======================================================================
\section{Nonlinear regime: step IC vs molecular dynamics}
\label{sec:nonlinear}
%======================================================================

We compare the DVM solution of Eq.~\eqref{eq:ghdb} from a
step velocity initial condition against the event-driven
molecular-dynamics trajectories of Ref.~\cite{PaperIII} on
two predictions of Ref.~\cite{PaperIV}: (a) the
linear-regime odd-mode amplitude ratio Eq.~\eqref{eq:ratio},
and (b) the even-mode saturation set by the collision
integral. The principal observables are the density-mode
Fourier amplitudes
$\hat\rho_n(t) = \tfrac{2}{N_s}\sum_k\delta_k(t)e^{-i n k_1 s_k}$
for $n=1,\dots,7$, sampled over twelve sound periods.

\subsection{Initial condition and parameters}

The step-velocity initial condition is
\begin{equation}
\bar u(s,0) = U_0\,\mathrm{sgn}\bigl[\cos(k_1 s)\bigr],
\quad
n_\alpha(s,0) = \tfrac12\rho^{\rm tot},
\label{eq:stepIC}
\end{equation}
with $k_1=2\pi/L$ and the velocity distribution at each $s$
the local Maxwell--Boltzmann Eq.~\eqref{eq:rho_eq} at the
common temperature $T=1$. We run two amplitudes:
$U_0=0.05$ (linear regime, where the $1/n$ ratio
Eq.~\eqref{eq:ratio} should be recovered) and $U_0=0.5$
(strongly nonlinear, where even modes become observable).
The simulation parameters are those of Sec.~\ref{sec:linear}
($L=12.57$, $\sigma=0.00314$, $\rho^{\rm tot}=31.82$,
$\eta=0.1$, $m_{\max}=5$, $N_s=80$, $N_v=128$,
$V_{\max}=6$, $\mathrm{CFL}=0.25$), integrated for
$T_{\rm sim}=12\,T_{\rm sound}\approx 135$ time units
($T_{\rm sound}=2\pi/(c_s^{\rm an} k_1)\approx 11.3$).

\subsection{Linear-regime mode ratio at $U_0=0.05$}

Figure~\ref{fig:peaks_lin} (data: \texttt{step\_sim\_step\_U005.csv})
shows the time evolution of $|\hat\rho_n(t)|$ for
$n=1,\dots,7$ over twelve sound periods at the linear
amplitude $U_0=0.05$. The peak amplitudes
$|\hat\rho_n|_{\max}\equiv\max_t|\hat\rho_n(t)|$ (after
discarding the initial $T_{\rm sound}/4$ to avoid the IC
transient) are reported in Table~\ref{tab:moderatio_lin}
together with the analytical $1/n$ prediction
Eq.~\eqref{eq:ratio}.

\begin{table}[t]
\caption{Linear-regime mode ratios at $U_0=0.05$
($m_{\max}=5$). DVM peaks measured at the FIRST local
maximum of $|\hat\rho_n(t)|$ (within the first sound period,
before higher-$n$ modes incur excess damping).}
\label{tab:moderatio_lin}
\centering
\begin{tabular}{c|ccc}
\hline\hline
$n$ & $|\hat\rho_n|^{\rm DVM}/|\hat\rho_1|^{\rm DVM}$
& $1/n$ & deviation (\%) \\
\hline
1 & $1.000$  & $1.000$ & $0.0$ \\
3 & $0.315$  & $0.333$ & $-5.6$ \\
5 & $0.180$  & $0.200$ & $-9.8$ \\
7 & $0.117$  & $0.143$ & $-18.1$ \\
\hline
2 (even) & $4.9\times 10^{-3}$ & $0$ & --- \\
4 (even) & $2.5\times 10^{-3}$ & $0$ & --- \\
6 (even) & $1.6\times 10^{-3}$ & $0$ & --- \\
\hline\hline
\end{tabular}
\end{table}

The odd-mode ratios recover the $1/n$ prediction of
Ref.~\cite{PaperIV} with a deviation that grows with $n$:
$5.6\%$ at $n=3$, $9.8\%$ at $n=5$, $18.1\%$ at $n=7$. The
deviation scales as $\sim n^2 k_1^2 D_{\rm app}\cdot T_{\rm sound}/2$,
identifying it as the integrated linear damping over a
fraction of the sound period that elapses between the
fundamental and the $n$-th harmonic crossing their respective
first peaks. The even modes at $U_0=0.05$ are at the
$0.5\%$ level relative to $|\hat\rho_1|$, consistent with
the expected $(U_0/c_s)^2\approx 0.2\%$ nonlinear coupling
plus a small discretisation-noise floor.

\subsection{Nonlinear even-mode saturation at $U_0=0.5$}

At $U_0=0.5\approx c_s^{\rm an}/2$ the dynamics enters the
nonlinear regime studied in Refs.~\cite{PaperIII,PaperIV}.
The principal qualitative result of Ref.~\cite{PaperIII} ---
absence of Burgers shock formation in the
mass-disordered hard-rod gas --- manifests in our DVM as a
SATURATION of the even-mode amplitudes at a finite fraction
of $|\hat\rho_1|$ (rather than the secular growth toward
$|\hat\rho_n|/|\hat\rho_1|\to 1$ characteristic of inviscid
Burgers steepening).

\begin{table}[t]
\caption{Nonlinear-regime mode ratios at $U_0=0.5$
($m_{\max}=5$, $U_0/c_s^{\rm an}=0.45$). FIRST-peak values.}
\label{tab:moderatio_nl}
\centering
\begin{tabular}{c|cc}
\hline\hline
$n$ & $|\hat\rho_n|^{\rm DVM}/|\hat\rho_1|^{\rm DVM}$
    & $1/n$ (linear ref) \\
\hline
1 & $1.000$ & $1.000$ \\
3 & $0.317$ & $0.333$ \\
5 & $0.177$ & $0.200$ \\
7 & $0.121$ & $0.143$ \\
\hline
2 (even) & $0.043$ & $0$ \\
4 (even) & $0.021$ & $0$ \\
6 (even) & $0.013$ & $0$ \\
\hline\hline
\end{tabular}
\end{table}

Table~\ref{tab:moderatio_nl} (data:
\texttt{step\_sim\_step\_U05.csv}) shows that the odd-mode
ratios are remarkably close to the linear-regime values of
Table~\ref{tab:moderatio_lin}, confirming the
linear-regime $1/n$ structure as the leading-order skeleton
of the nonlinear evolution: the ten-fold amplitude increase
$U_0:0.05\to 0.5$ leaves the odd ratios within $1\%$ of their
linear values (e.g.\ $0.315\to 0.317$ at $n=3$). The
even-mode amplitudes are no longer at the noise floor: they
acquire finite ratios
$|\hat\rho_{2,4,6}|/|\hat\rho_1|=0.043,0.021,0.013$,
i.e.\ $O(U_0/c_s)\approx 0.45$ scaling for $n=2$ and a
hierarchical decay $\sim 0.5\cdot$ per even step
that does NOT reach the Burgers shock fixed point
$|\hat\rho_n|/|\hat\rho_1|\to 1$. This is the saturation
mechanism proposed in Ref.~\cite{PaperIV} Sec.~VI.B,
predicted there to follow from collisional inhibition of
high-$n$ mode growth.

\subsection{Sound-mode damping rate}

The damping rate of the fundamental mode $|\hat\rho_1|$ is
extracted by linear regression of
$\log|\hat\rho_1|_{\rm peak}(t_k)$ against the times $t_k$ of
successive local maxima of $|\hat\rho_1(t)|$ over twelve
sound periods at the linear amplitude $U_0=0.05$. The fitted
value at $m_{\max}=5$ is
\begin{equation}
\gamma^{\rm DVM} = 4.23\times 10^{-2},
\quad
Q^{\rm DVM} = \frac{\omega_1}{2\gamma^{\rm DVM}} = 6.6,
\end{equation}
yielding an apparent diffusion coefficient
$D_{\rm app}=\gamma^{\rm DVM}/k_1^2=0.17$. The
Chapman--Enskog prediction of Sec.~\ref{sec:nu} is
$\gamma_{\rm CE}\sim\kappa D_{LH}k_1^2\approx 0.008$ for
$\kappa\sim 1/(4 \rho^{\rm tot})$; the empirical
$\gamma_{\rm MD}\sim 0.02$--$\,0.03$ of Ref.~\cite{PaperII}
falls between these two values. The DVM damping
overestimates the empirical $\gamma_{\rm MD}$ by a factor
$\sim 1.7$, which we attribute to the $O(\Delta v^2)$
numerical viscosity of the bilinear collision kernel
(Sec.~\ref{sec:dvm}); subtracting a baseline $\gamma_{\rm num}$
estimated from a collisionless control run would close this
gap and is left to the conservative-kernel upgrade flagged
in Sec.~\ref{sec:discussion}. The quality factor $Q\sim 7$
is moderately underdamped, consistent with the
$Q_{\rm MD}\sim 11$--$\,16$ reported in Ref.~\cite{PaperII}
to within the same numerical-viscosity factor.

\subsection{Limits of the present comparison}

Two systematic effects limit the precision of the DVM--MD
comparison at the present implementation. First, the
$O(\Delta v^2)$ conservation error of the bilinear collision
kernel introduces a small numerical contribution to the
damping rate ($\gamma_{\rm num}$), which acts as an
additive baseline that is difficult to separate from the
genuine $\gamma_{\rm CE}$ at $N_v=128$; a clean separation
would require either the 3-pt Lagrange kernel of
Sec.~\ref{sec:nu} ported to the full time-stepping code, or
the explicit subtraction of $\gamma_{\rm num}$ measured in
a collisionless control run. Second, the comparison against
Ref.~\cite{PaperIII} is at present qualitative (peak ratios
and saturation behaviour) rather than trajectory-resolved,
because Ref.~\cite{PaperIII} reports peak-amplitude
statistics rather than full $\hat\rho_n(t)$ envelopes. A
trajectory-resolved comparison would require re-running the
MD with matched output cadence and is left to a companion
data release.

%======================================================================
\section{Discussion}\label{sec:discussion}
%======================================================================

\subsection{Closure of the open items of Ref.~\cite{PaperIV}}

The construction of this paper closes three of the four open
quantitative items of Ref.~\cite{PaperIV}:

\emph{(i) Sound-wave damping rate.} Identified in
Sec.~\ref{sec:nu} as the inter-species interdiffusion
$D_{LH}(m_{\max})$ rather than a conventional shear
viscosity, and computed from first principles via the
Chapman--Enskog expansion of the linearised GHD--Boltzmann
operator with the 3-point Lagrange conservative kernel of
Sec.~\ref{sec:nu}. The numerical extraction is essentially
flat in $m_{\max}$, $D_{LH}\approx 4.0$ at the parameters of
Ref.~\cite{PaperII}, and yields a sound-mode damping rate
in agreement with the empirical
$\gamma_{\rm MD}\sim 0.02$--$\,0.03$ of Ref.~\cite{PaperII}
to within a factor of $\sim 3$ pending the
thermodynamic-prefactor $\kappa(m_L,m_H,T,n)$ derivation
that is left to future work.

\emph{(ii) Even-mode saturation.} Sec.~\ref{sec:nonlinear}
shows that the DVM solver reproduces the saturation
qualitatively at $U_0=0.5$: the even modes acquire finite
amplitudes saturating at
$|\hat\rho_2|/|\hat\rho_1|\sim U_0/c_s$ rather than growing
secularly toward the inviscid-Burgers shock limit of
Ref.~\cite{PaperIII}.

\emph{(iv) Direct nonlinear-regime test.} The DVM matches
the linear-regime $1/n$ ratio of Ref.~\cite{PaperIV} at
$U_0=0.05$ and the even-mode saturation of
Ref.~\cite{PaperIII} at $U_0=0.5$ within the precision of
the bilinear collision kernel; a quantitatively tight
trajectory-resolved comparison awaits both the conservative-%
kernel upgrade and the matched-cadence MD reproduction noted
in Sec.~\ref{sec:nonlinear}.

The remaining item \emph{(iii)} --- the microscopic
derivation of the dressed-velocity postulate
Eq.~\eqref{eq:veff} --- is unresolved by the present
numerical approach and requires the
unequal-mass Jepsen-mapping extension flagged in
Ref.~\cite{PaperIV} Sec.~VII.D.

\subsection{Limits of the DVM approach}

Three principal limits of the present implementation should
be noted. First, the streaming substep uses a WENO5
upwind reconstruction with $O(\Delta s^5)$ accuracy in
smooth regions, which is sufficient for the sound-speed
recovery within $0.1\%$ at $N_s=80$ reported in
Sec.~\ref{sec:linear}. The choice of fit window in the
sound-speed extraction is essential at this precision
level: a small left-moving component of the
initial-condition eigenmode introduces a phase modulation
that, combined with the damping rate, biases long-window
fits toward smaller apparent $c_s$
(Sec.~\ref{sec:linear}).
Second, the dressed-velocity streaming Eq.~\eqref{eq:veff} is
singular as $\eta\to 1$ and we clamp $\eta\leq 0.99$ in our
implementation; this is never activated at our chosen
$\eta=0.1$ but limits the applicability of the present
solver to the dilute and moderate-density regimes where
$\eta\lesssim 0.5$.
Third, the velocity grid must contain the relevant thermal
and coherent scales for both species; at $m_{\max}=100$ the
heavy-species thermal width is $v_{\rm th}^H=0.1$ and
$V_{\max}/v_{\rm th}^H=60$ at $V_{\max}=6$, leaving only
$\sim 2\,N_v/60\approx 4$ velocity bins per heavy-thermal
width, which is marginal for high-amplitude transients and
should be improved by adaptive velocity gridding in a
follow-up.

\subsection{Connection to the P--V/P--VI sequence}

The numerical machinery developed here feeds directly into
the traffic application of the companion proceedings paper
\cite{PaperVI}: the GHD--Boltzmann framework
generalised in this paper to interpolated and conservative
discrete-velocity simulation provides the kinetic
foundation for the inelastic and delay-kernel extensions of
the hard-rod gas required for traffic modelling. The
3-point Lagrange kernel of Sec.~\ref{sec:nu} is the precise
discretisation that the inelastic kernel of
Ref.~\cite{PaperVI} App.~A requires for its discrete-energy
balance; the C--E interdiffusion-coefficient framework of
Sec.~\ref{sec:nu} is the precise transport-coefficient
framework that the traffic threshold-prediction analysis of
Ref.~\cite{PaperVI} Sec.~V uses for its effective-shear-%
viscosity estimate.

A natural next step is a trajectory-resolved DVM--MD
comparison at matched output cadence, alongside a
recalculation of $\gamma$ with subtracted numerical-damping
baseline (collisionless control run) to disentangle the
genuine $\gamma_{\rm CE}$ of Sec.~\ref{sec:nu} from the
residual streaming-induced numerical dissipation of the
WENO5 scheme on the present $N_s=80$ grid.

\begin{acknowledgments}
This work was supported by Tecnol\'ogico de Monterrey.
Computations were performed on the xolotl cluster. The Julia
implementation of the GHD--Boltzmann solver and the data of
the reported runs will be made publicly available upon
publication.
\end{acknowledgments}

\begin{thebibliography}{99}

\bibitem{PaperI}
H.~Medel, in preparation (2026).

\bibitem{PaperII}
H.~Medel, in preparation (2026).

\bibitem{PaperIII}
H.~Medel, in preparation (2026).

\bibitem{PaperIV}
H.~Medel, in preparation (2026).

\bibitem{DoyonSpohn2017}
B.~Doyon, H.~Spohn,
J.~Stat.~Mech.\ 073210 (2017).

\bibitem{GHDperspective2025}
B.~Doyon, S.~Gopalakrishnan, F.~M{\o}ller,
J.~Schmiedmayer, R.~Vasseur,
Phys.~Rev.~X \textbf{15}, 010501 (2025).

\bibitem{HubnerDoyon2024}
F.~H\"ubner, B.~Doyon,
arXiv:2406.18322 (2024).

\bibitem{HarrisHurst}
H.~Spohn,
\emph{Large-Scale Dynamics of Interacting Particles},
Springer (1991);
J.~L.~Lebowitz, J.~Percus, J.~Sykes,
Phys.~Rev.\ \textbf{171}, 224 (1968).

\bibitem{ChapmanCowling}
S.~Chapman and T.~G.~Cowling,
\emph{The Mathematical Theory of Non-Uniform Gases},
3rd ed., Cambridge University Press (1970).

\bibitem{Mieussens2000}
L.~Mieussens,
\emph{Discrete-velocity models and numerical schemes for the
Boltzmann--BGK equation in plane and axisymmetric geometries},
J.~Comput.~Phys.\ \textbf{162}, 429 (2000).

\bibitem{PaperVI}
H.~Medel, in preparation (2026).

\end{thebibliography}

\end{document}
